% !TEX root = ../auv-thesis.tex
% Canonical source: docs/thesis/paper/06_conclusion_and_outlook.md
% Source SHA256: ff6d6c0e0c539007b90b73c417c8648ba120b9c50abee3d8ae04d835642e3276
% Generated by tools/migrate_thesis_chapters.py; do not edit independently.

\chapter{结论与展望}

\section{全文总结}\label{sec:ch06-6-1}

本文围绕``AUV 自主海缆巡检''这一工程命题，针对深远海海缆资产快速增长而传统巡检手段在成本、效率与安全性上难以为继的现实矛盾，提出并实现了一套面向不确定性、以协方差流转为主线的自主巡检软件系统，并在虚实结合的实验体系中对其算法层面的正确性与边界给出了可追溯的证据。全文的核心工作可归纳为以下四个方面。

第一，在系统层面，本文确立了感知、估计、决策与控制四环贯通，不确定性沿层向上传递的整体架构，并将其落地为 Jetson Orin 与 PC104 双脑异步协同的软硬件形态。与把各功能模块简单串接的常规做法不同，本文的立足点在于让每一层不仅向上层输出状态量，还输出对该状态量的置信度（协方差），使上层决策与控制能够``知道自己知道多少''。这一设计把``不确定性感知''从一句口号转化为可在模块接口上核查的数据契约，也构成了后续状态估计与控制章节的共同基础。

第二，在状态估计层面，本文以误差状态扩展卡尔曼滤波（ES-EKF）为骨架，融合惯性、DVL、声学与磁场多源观测，构建了声磁协同的近底定位与海缆路由跟踪方案，并引入基于归一化新息平方（NIS）的自适应机制在线调节量测噪声。实验表明，该方案在八类传感器与通信扰动场景、三种子共二十四次运行下全部成功，且协方差一致性聚合显示自适应机制在全场景被触发，说明状态估计在多类扰动下既能保持稳定、其协方差又能保持可核查的一致性，这正是把置信度向控制层可靠传递的前提。

第三，在决策与控制层面，本文采用行为树组织任务级决策、以分层任务状态机与应急切换逻辑保证安全兜底，并在运动控制层实现了将在线感知置信度纳入代价权重调制的不确定性感知模型预测控制（UA-MPC）。公平口径下的对照实验表明，模型预测控制在长短波 S 弯与发卡掉头等不规则曲率路径上优于或持平于视线制导与比例积分微分基线，其优势区恰是几何预瞄需求最强之处；不确定性感知控制在基线与综合压力场景下展现出定位与控制双侧的改善（水平误差约 10.4\%、横向误差约 15.8\%），但在重度 DVL 丢包场景下不再具备优势。这一``有条件改善''的结果，恰恰印证了``感知与控制的不确定性必须协同处理、单一环节无法包打天下''的整体判断。

第四，在实验验证层面，本文构建了数值分析、软件在环、硬件在环与实验室四层虚实结合体系，并按证据等级对每条结论标注了样本量与适用边界。声磁电缆探测的端到端链路已在干净先验下闭环并通过行业标准数字孪生验收，在先验畸变下于六自由度闭环中首次复现了闭环恢复；磁传感器杆臂标定完成了一轮仿真标定验证（平移误差降低 96.27\%、旋转降低 95.31\%）。硬件文档证据方面，编号 6\# TMR8637 随附三轴测试报告给出 1 Hz 处 200 pT/\(\sqrt{\text{Hz}}\) 和 50 Hz 附近约 20--30 pT/\(\sqrt{\text{Hz}}\) 的模组本底噪声，说明模组本体具备支撑 0.05 nT 窄带系统目标的噪声裕度，但完整采集链仍待实测验收。这些结果共同支撑了``算法层面正确性、模组本体证据加已识别系统硬件接口''的分层证据图。

综合来看，本文的主要贡献不在于宣称某一算法全面领先，而在于提出并验证了一条``让不确定性在系统各层显式流转、并据此调制决策与控制''的技术路线，并以分层、可核查的实验证据界定了这条路线当前的能力边界。

\section{研究局限}\label{sec:ch06-6-2}

严守证据等级是本文的一贯立场，因此有必要把当前工作尚不能支撑的结论明确列出，避免读者做出超出证据范围的推断。

其一，实验证据整体仍处于仿真与半实物阶段，系统级硬件物理证据仍不完整。TMR8637 随附测试报告只补上单件模组的三轴传输与本底噪声证据，不能替代 SK2301、隔离供电、地环路、车体自生磁场和 DLIA 组成的完整链路验收。磁传感器九参数标定、完整链路底噪谱、实测埋深反演与真机双脑时延目前均属``实验方案加未来工作''；仿真平台时延不能等同于真机时延，仿真标定框架也不能等同于真机转台标定。

其二，部分结论的统计充分性不足。不确定性感知控制的主消融、极端电缆场景的代理冒烟测试等多为单种子或少种子结果，只能写成``可运行的压力测试与待统计趋势''，不能上升为固定权重与不确定性感知之间的统计优劣结论。

其三，场景真实性存在差距。当前仿真尚未原生支持``电缆几何、近底地形、声磁观测与横流''四者耦合的完整数字孪生，极端场景通过代理方式表达，与真实海试之间仍有明确差距；回退路径在仿真侧也尚未被极端场景有效压力测试。

其四，可交付的实物部署验收仍为否。满足磁观测前提的六自由度闭环恢复虽已复现，但真实检测噪声、多种子统计与硬件实物三个环节尚未闭合，因此不能表述为``通过真实海缆检测精度验收''。

这些局限并非本文路线的反例，而是从仿真验证走向工程交付过程中自然存在的待办项，它们直接构成了下一节展望的出发点。

\section{未来工作展望}\label{sec:ch06-6-3}

结合上述局限，未来工作可沿四个方向递进推进，前两项属短中期的仿真内补强，后两项属面向工程交付与能力扩展的长期方向。

\textbf{其一，缩小仿真与现实的差距，推进 Sim-to-Real 迁移。} 当前证据的最大短板在于场景真实性与系统级硬件验证不足，因此首要工作是把算法从数字孪生稳健地迁移到真机。域随机化通过在训练阶段主动扰动仿真参数以覆盖真实分布，被证明是弥合仿真与现实差距的有效途径\cite{tobin2017domain}；进一步地，以数字孪生为闭环载体、驱动 AUV 集群协同探测的最新实践，为仿真训练、真机部署与在线回灌的持续迭代提供了可借鉴的范式\cite{yan2026digitaltwin}。后续应在本文的四层实验体系上补齐真实检测噪声注入与多种子统计，实测 TMR8637-\/-SK2301-\/-DLIA 完整链路的 50 Hz 输出噪声与等效噪声带宽，并逐步引入真机转台标定与实海况数据回灌。

\textbf{其二，深化不确定性感知控制，引入学习型与随机模型预测控制。} 本文的 UA-MPC 以人工设计的置信度与权重映射为主，未来可向数据驱动的自适应方向演进。学习型模型预测控制系统性地讨论了如何在保证安全的前提下用数据改进模型与代价，为在线学习环境模型提供了理论框架\cite{hewing2020learning}；带不确定性传播时域的随机非线性模型预测控制，则从概率约束的角度给出了在预测时域内显式传播不确定性的具体做法，与本文``协方差沿层流转''的主线高度契合，可作为把置信度更严格地纳入优化问题的下一步\cite{zarrouki2024stochastic}。

\textbf{其三，提升海缆探测的智能化水平。} 本文的探测链路以物理建模与解析反演为主，未来可引入面向 AUV 巡检的实时智能检测方法，把学习型感知与本文的声磁协同估计相结合，以提升在复杂地形与强干扰下对海缆目标的实时识别与跟踪能力\cite{bodenmann2024realtime}，从而把``前置筛查''的召回率与鲁棒性进一步抬高。

\textbf{其四，扩展作业能力，迈向多 AUV 协同与持久自治。} 面向深远海大规模周期性巡检，单台 AUV 的续航与覆盖能力终将成为瓶颈。一方面可发展多 AUV 协同编队以并行覆盖长距离路由。多自主水下航行器编队控制的系统综述表明，领航跟随、虚拟结构与基于行为等编队范式各有适用边界，而水声通信的低带宽、高时延与拓扑时变则是其区别于陆空编队的核心约束\cite{yan2023formation}；与之相辅相成的是协同定位，即让编队内的 AUV 通过相互测距与信息交换来抑制各自惯导的长期漂移。近期基于最大熵原理的鲁棒协同定位方法进一步说明了在量测野值与非高斯噪声下如何保持定位一致性\cite{li2025cooplocalization}，这与本文``协方差沿层流转、并以一致性检验守护估计可信度''的主线在多机层面上完全同构，可视为把单机不确定性感知思路推广到编队尺度的自然延伸。另一方面，驻留式 AUV 与水下对接技术为长航时、免母船值守的持久自治作业提供了可行路径\cite{curtin2018resident}，使 AUV 能够在海底基站就地充电与数据回传，从而把巡检从``单次出海任务''升级为``常驻式持续监测''。

总体而言，上述四个方向由近及远地勾勒出从``算法层面正确''走向``工程可交付''的完整路线：短中期以仿真内的统计补强与场景扩展为主，长期以硬件实物验证与多机协同持久自治为目标，二者按``先仿真扩展、再迁到真机''的顺序推进，即可在不依赖一次性满足所有条件的前提下，逐步把本文技术路线的实证强度抬升到工程交付水平。
